\chapter{Differential forms}
\label{chap:forms}

Differential forms on manifolds modeled on locally convex spaces and the exterior derivative,
following \cite[Section~3.6]{GN26}. Two features of the finite-dimensional theory are missing:
there is no coordinate description of a form by components in a chart, and there are no
partitions of unity, even on Banach manifolds. The definitions below only use local
representatives, and all statements are local. No completeness of the model space or of the
space of values is assumed.

Glöckner and Neeb define the exterior derivative through the Lie algebra differential of vector
fields \cite[Proposition~3.6.7, formula~(3.29)]{GN26} and show that it is given in charts by the
formula~(3.33) below. We take~(3.33) as the definition and check directly that it does not depend
on the chart: the terms with second derivatives of the transition maps cancel because these are
symmetric. The same kind of computation proves $\dd \circ \dd = 0$, as indicated in
\cite[Remark~3.6.10]{GN26}.


\section{Local forms}

Let $E$ and $F$ be Hausdorff locally convex spaces and $U \subseteq E$ open.

\begin{definition}[Local forms]
  \label{def:local-form}
  \lean{LocallyConvex.Manifolds.IsLocalForm}
  \leanok
  A \emph{local $F$-valued $n$-form} on $U$ is a map $\omega \colon U \times E^n \to F$ which is
  smooth in the sense of Michal--Bastiani and alternating $n$-linear in the last $n$ variables.
\end{definition}

\begin{definition}[Local exterior derivative]
  \label{def:local-d}
  \lean{LocallyConvex.Manifolds.dLoc}
  \leanok
  \uses{def:local-form}
  For a local $n$-form $\omega$,
  \[
    \dd\omega(y)(v_0, \dots, v_n) =
      \sum_{i=0}^{n} (-1)^i \, \partial_{v_i} \omega(y)(v_0, \dots, \widehat{v_i}, \dots, v_n),
  \]
  where $\partial_{v_i}$ is the directional derivative in the base point $y$
  \cite[(3.33)]{GN26}.
\end{definition}

\begin{proposition}
  \label{prop:local-d-form}
  \lean{LocallyConvex.Manifolds.IsLocalForm.dLoc, LocallyConvex.Manifolds.IsLocalForm.exists_partial}
  \leanok
  \uses{def:local-d}
  The exterior derivative of a local $n$-form is a local $(n+1)$-form.
\end{proposition}

\begin{proof}
  \leanok
  Smoothness: each term is the derivative $\dd\omega$ composed with a continuous linear map.
  For fixed $y$, the partial derivative $(u, w) \mapsto \partial_u \omega(y)(w)$ is linear in $u$
  and alternating $n$-linear in $w$, because $\omega(z)$ is so for every $z$ near $y$; then
  $\dd\omega(y)$ is its alternation, which is alternating $(n+1)$-linear.
\end{proof}

\begin{theorem}[$\dd \circ \dd = 0$, locally]
  \label{thm:local-dd}
  \lean{LocallyConvex.Manifolds.IsLocalForm.dLoc_dLoc,
    LocallyConvex.Manifolds.sum_neg_one_pow_smul_sum_eq_zero_of_symm}
  \leanok
  \uses{prop:local-d-form}
  For a local form $\omega$ on $U$, $\dd(\dd\omega) = 0$ on $U$.
\end{theorem}

\begin{proof}
  \leanok
  $\dd(\dd\omega)(y)(v_0, \dots, v_{n+1})$ is a double alternating sum of second derivatives
  $\partial_{v_i}\partial_{v_j}\omega(y)$, which are symmetric in $i, j$ (Schwarz's theorem for
  Michal--Bastiani maps). The terms cancel in pairs \cite[Remark~3.6.10]{GN26}.
\end{proof}

\begin{definition}[Local pullback]
  \label{def:local-pullback}
  \lean{LocallyConvex.Manifolds.pullbackLoc, LocallyConvex.Manifolds.IsLocalForm.pullback}
  \leanok
  \uses{def:local-form}
  For a smooth map $\tau \colon U \to V$ between open subsets of locally convex spaces and a
  local form $\omega$ on $V$,
  $(\tau^*\omega)(y)(v_1, \dots, v_n) = \omega(\tau y)(\dd\tau(y) v_1, \dots, \dd\tau(y) v_n)$.
  It is a local form on $U$.
\end{definition}

\begin{theorem}[Naturality, locally]
  \label{thm:local-naturality}
  \lean{LocallyConvex.Manifolds.IsLocalForm.dLoc_pullbackLoc,
    LocallyConvex.Manifolds.IsLocalForm.dMB_pullbackLoc,
    AlternatingMap.sum_neg_one_pow_smul_sum_update_eq_zero}
  \leanok
  \uses{def:local-pullback, def:local-d}
  $\dd(\tau^*\omega) = \tau^*(\dd\omega)$ on $U$.
\end{theorem}

\begin{proof}
  \leanok
  By the chain rule, $\partial_u (\tau^*\omega)(y)(w)$ is
  $\partial_{\dd\tau(y) u}\omega(\tau y)(\dd\tau(y) w)$ plus the sum over $k$ of
  $\omega(\tau y)(\dd\tau(y) w_1, \dots, \dd^2\tau(y)(w_k, u), \dots, \dd\tau(y) w_n)$; the second
  part uses that $\omega(\tau y)$ is multilinear. In the alternating sum, the first terms give
  $\tau^*(\dd\omega)$, and the terms with $\dd^2\tau(y)$ cancel because $\dd^2\tau(y)$ is
  symmetric.
\end{proof}

\section{Forms on manifolds}

Let $M$ be a smooth manifold modeled on $E$. A tangent vector at $x$ is given by its coordinates
in the preferred chart at $x$ (Definition~\ref{def:tangent-bundle}).

\begin{definition}[Differential forms]
  \label{def:diff-form}
  \lean{LocallyConvex.Manifolds.DifferentialForm, LocallyConvex.Manifolds.formLocalRep,
    LocallyConvex.Manifolds.DifferentialForm.localRep}
  \leanok
  \uses{def:tangent-bundle, def:local-form}
  An \emph{$F$-valued $n$-form} $\omega$ on $M$ assigns to each $x \in M$ an alternating
  $n$-linear map $\omega_x \colon (T_xM)^n \to F$ such that, for every chart $\varphi$ of the
  atlas, the local representative
  $(y, v_1, \dots, v_n) \mapsto \omega_{\varphi^{-1} y}(T\varphi^{-1} v_1, \dots,
  T\varphi^{-1} v_n)$ is smooth on $\varphi(U) \times E^n$ \cite[Definition~3.6.1]{GN26}. We
  write $\Omega^n(M, F)$ for the space of these forms.
\end{definition}

\begin{lemma}
  \label{lem:form-localrep}
  \lean{LocallyConvex.Manifolds.DifferentialForm.isLocalForm_localRep,
    LocallyConvex.Manifolds.DifferentialForm.localRep_eq_pullbackLoc,
    LocallyConvex.Manifolds.DifferentialForm.localRep_prefChart}
  \leanok
  \uses{def:diff-form, def:local-pullback}
  The local representatives of a form are local forms. On chart overlaps they are related by
  pullback along the transition map: $\omega_i = \tau^*\omega_j$ for $\tau = j \circ i^{-1}$.
\end{lemma}

\begin{proof}
  \leanok
  The coordinate changes of tangent vectors are the derivatives of the transition maps, and they
  compose by the chain rule.
\end{proof}

\begin{definition}[Exterior derivative]
  \label{def:diff-form-d}
  \lean{LocallyConvex.Manifolds.DifferentialForm.d, LocallyConvex.Manifolds.DifferentialForm.d_apply}
  \leanok
  \uses{def:diff-form, def:local-d}
  For $\omega \in \Omega^n(M, F)$, $(\dd\omega)_x$ is the local exterior derivative of the local
  representative of $\omega$ in the preferred chart at $x$, evaluated at the image of $x$.
\end{definition}

\begin{theorem}
  \label{thm:diff-form-d}
  \lean{LocallyConvex.Manifolds.DifferentialForm.localRep_d,
    LocallyConvex.Manifolds.DifferentialForm.formLocalRep_dFun}
  \leanok
  \uses{def:diff-form-d, lem:form-localrep, thm:local-naturality, prop:local-d-form}
  In every chart of the atlas, the local representative of $\dd\omega$ is the local exterior
  derivative of the local representative of $\omega$. In particular, $\dd\omega$ is a smooth
  $(n+1)$-form \cite[Proposition~3.6.7]{GN26}.
\end{theorem}

\begin{proof}
  \leanok
  Near a point, $\omega_i = \tau^*\omega_j$ with $\tau$ the transition map from the chart $i$ to
  the preferred chart $j$ at the point; apply Theorem~\ref{thm:local-naturality}.
\end{proof}

\begin{theorem}[$\dd \circ \dd = 0$]
  \label{thm:diff-form-dd}
  \lean{LocallyConvex.Manifolds.DifferentialForm.d_d}
  \leanok
  \uses{thm:diff-form-d, thm:local-dd}
  For every $\omega \in \Omega^n(M, F)$, $\dd(\dd\omega) = 0$ \cite[Proposition~3.6.9]{GN26}.
\end{theorem}

\begin{proof}
  \leanok
  In the preferred chart at $x$, the local representative of $\dd\omega$ is $\dd$ of the local
  representative of $\omega$ (Theorem~\ref{thm:diff-form-d}); apply
  Theorem~\ref{thm:local-dd}.
\end{proof}

\begin{definition}[Functions, closed and exact forms]
  \label{def:closed-exact}
  \lean{LocallyConvex.Manifolds.DifferentialForm.ofFunction,
    LocallyConvex.Manifolds.DifferentialForm.IsClosed,
    LocallyConvex.Manifolds.DifferentialForm.IsExact,
    LocallyConvex.Manifolds.DifferentialForm.IsExact.isClosed}
  \leanok
  \uses{def:diff-form-d, thm:diff-form-dd}
  A smooth function $f \colon M \to F$ is an $F$-valued $0$-form. A form $\omega$ is
  \emph{closed} if $\dd\omega = 0$ and \emph{exact} if $\omega = \dd\eta$ for some form $\eta$
  \cite[Definition~3.6.11]{GN26}. By Theorem~\ref{thm:diff-form-dd}, exact forms are closed.
\end{definition}

\begin{proposition}
  \label{prop:d-function}
  \lean{LocallyConvex.Manifolds.DifferentialForm.d_ofFunction_apply}
  \leanok
  \uses{def:closed-exact}
  For a smooth function $f$, $(\dd f)_x(v) = \dd(f \circ \varphi^{-1})(\varphi x)\, v$, where
  $\varphi$ is the preferred chart at $x$: the exterior derivative of $f$ is its differential
  \cite[Remark~3.6.12(a)]{GN26}.
\end{proposition}

\begin{proof}
  \leanok
  For $n = 0$ the alternating sum of Definition~\ref{def:local-d} has a single term.
\end{proof}

\section{The vector space of forms}

\begin{definition}[Linear structure]
  \label{def:forms-linear}
  \lean{LocallyConvex.Manifolds.DifferentialForm.instModule,
    LocallyConvex.Manifolds.DifferentialForm.add_apply,
    LocallyConvex.Manifolds.DifferentialForm.smul_apply}
  \leanok
  \uses{def:diff-form}
  Sums and real multiples of forms are taken pointwise; their local representatives are the
  corresponding combinations of local representatives, hence smooth. This makes
  $\Omega^n(M, F)$ a real vector space.
\end{definition}

\begin{proposition}
  \label{prop:d-linear}
  \lean{LocallyConvex.Manifolds.DifferentialForm.dLinear,
    LocallyConvex.Manifolds.DifferentialForm.d_add,
    LocallyConvex.Manifolds.DifferentialForm.d_smul}
  \leanok
  \uses{def:forms-linear, def:diff-form-d}
  The exterior derivative $\dd \colon \Omega^n(M, F) \to \Omega^{n+1}(M, F)$ is linear.
\end{proposition}

\begin{proof}
  \leanok
  In the preferred chart at each point, directional derivatives are additive and homogeneous in
  the function.
\end{proof}

\section{Pullbacks}

\begin{definition}[Pullback]
  \label{def:diff-form-pullback}
  \lean{LocallyConvex.Manifolds.DifferentialForm.pullback,
    LocallyConvex.Manifolds.DifferentialForm.formLocalRep_pullbackFun}
  \leanok
  \uses{def:diff-form, def:tangent-map, def:local-pullback}
  For a smooth map $\varphi \colon N \to M$ and $\omega \in \Omega^n(M, F)$,
  $(\varphi^*\omega)_x(v_1, \dots, v_n) = \omega_{\varphi(x)}(T_x\varphi\, v_1, \dots,
  T_x\varphi\, v_n)$ \cite[Definition~3.6.6]{GN26}. In charts, the local representative of
  $\varphi^*\omega$ is the local pullback of the local representative of $\omega$ along the local
  representative of $\varphi$; in particular $\varphi^*\omega$ is a smooth form.
\end{definition}

\begin{theorem}[Naturality of $\dd$]
  \label{thm:diff-form-naturality}
  \lean{LocallyConvex.Manifolds.DifferentialForm.d_pullback}
  \leanok
  \uses{def:diff-form-pullback, thm:diff-form-d, thm:local-naturality}
  For a smooth map $\varphi \colon N \to M$ and $\omega \in \Omega^n(M, F)$,
  $\dd(\varphi^*\omega) = \varphi^*(\dd\omega)$ \cite[Lemma~3.6.13]{GN26}.
\end{theorem}

\begin{proof}
  \leanok
  In the preferred charts at $x$ and $\varphi(x)$, this is Theorem~\ref{thm:local-naturality}
  for the local representative of $\varphi$.
\end{proof}

\begin{proposition}[Functoriality of pullbacks]
  \label{prop:pullback-functorial}
  \lean{LocallyConvex.Manifolds.DifferentialForm.pullback_id,
    LocallyConvex.Manifolds.DifferentialForm.pullback_pullback,
    LocallyConvex.Manifolds.DifferentialForm.pullbackLinear}
  \leanok
  \uses{def:diff-form-pullback, thm:man-tangent-map, def:forms-linear}
  The pullback $\varphi^*$ is linear, $\id^*\omega = \omega$ and
  $\psi^*(\varphi^*\omega) = (\varphi \circ \psi)^*\omega$ \cite[(3.27)]{GN26}.
\end{proposition}

\begin{proof}
  \leanok
  The vectors in $\varphi^*\omega$ are images under the tangent map, so this is the chain rule
  $T(\varphi \circ \psi) = T\varphi \circ T\psi$ and $T\id = \id$.
\end{proof}

\begin{proposition}[Composition with linear maps]
  \label{prop:d-comp-clm}
  \lean{LocallyConvex.Manifolds.DifferentialForm.compCLM,
    LocallyConvex.Manifolds.DifferentialForm.d_compCLM}
  \leanok
  \uses{def:diff-form-d}
  For a continuous linear map $\alpha \colon F \to G$, $(\alpha_*\omega)_x = \alpha \circ
  \omega_x$ defines a $G$-valued form, and $\alpha_*(\dd\omega) = \dd(\alpha_*\omega)$
  \cite[Exercise~3.6.4]{GN26}.
\end{proposition}

\begin{proof}
  \leanok
  A continuous linear map commutes with directional derivatives.
\end{proof}

\section{The wedge product}

Let $\beta \colon E_1 \times E_2 \to E_3$ be a continuous bilinear map between locally convex
spaces.

\begin{definition}[Wedge product]
  \label{def:wedge}
  \lean{LocallyConvex.Manifolds.wedgeFun, LocallyConvex.Manifolds.wedgeAlt,
    LocallyConvex.Manifolds.wedgeAlt_apply, LocallyConvex.Manifolds.DifferentialForm.wedge}
  \leanok
  \uses{def:diff-form}
  For $\omega \in \Omega^p(M, E_1)$ and $\eta \in \Omega^q(M, E_2)$,
  \[
    (\omega \wedge \eta)_x(v_1, \dots, v_{p+q}) = \frac{1}{p!\,q!}
      \sum_{\sigma \in S_{p+q}} \operatorname{sgn}(\sigma)\,
      \beta\bigl(\omega_x(v_{\sigma(1)}, \dots, v_{\sigma(p)}),
        \eta_x(v_{\sigma(p+1)}, \dots, v_{\sigma(p+q)})\bigr)
  \]
  \cite[Definition~3.6.3]{GN26}. It is alternating, being the alternation of a multilinear map,
  and its local representatives are finite sums of compositions of $\beta$ with local
  representatives of $\omega$ and $\eta$, hence smooth.
\end{definition}

\begin{proposition}
  \label{prop:wedge-properties}
  \lean{LocallyConvex.Manifolds.DifferentialForm.wedgeBilinear,
    LocallyConvex.Manifolds.DifferentialForm.wedge_apply_one_one,
    LocallyConvex.Manifolds.DifferentialForm.pullback_wedge}
  \leanok
  \uses{def:wedge, def:diff-form-pullback, def:forms-linear}
  The wedge product is bilinear; for $p = q = 1$,
  $(\omega \wedge \eta)_x(v_1, v_2) = \beta(\omega_x v_1, \eta_x v_2) - \beta(\omega_x v_2,
  \eta_x v_1)$ \cite[Definition~3.6.3]{GN26}; and $\varphi^*(\omega \wedge \eta) =
  \varphi^*\omega \wedge \varphi^*\eta$ \cite[(3.28)]{GN26}.
\end{proposition}

\begin{proof}
  \leanok
  Direct from the definitions; for $p = q = 1$ the symmetric group has two elements.
\end{proof}

\begin{proposition}[Products with functions]
  \label{prop:smul-fun}
  \lean{LocallyConvex.Manifolds.DifferentialForm.smulFun,
    LocallyConvex.Manifolds.DifferentialForm.mul_smulFun,
    LocallyConvex.Manifolds.DifferentialForm.wedge_ofFunction}
  \leanok
  \uses{def:wedge, def:closed-exact}
  For a smooth function $f \colon M \to \R$, $(f\omega)_x = f(x)\,\omega_x$ defines a smooth form,
  and these products satisfy the module identities over $C^\infty(M, \R)$. The form $f\omega$ is
  the wedge product of $\omega$ with the $0$-form $f$ for the scalar multiplication
  \cite[Remark~3.6.5]{GN26}.
\end{proposition}

\begin{proof}
  \leanok
  Smoothness follows from the product rule for the scalar multiplication $\R \times F \to F$.
  For the wedge product with a $0$-form, each permutation contributes $\operatorname{sgn}(\sigma)^2
  f(x)\,\omega_x(v) = f(x)\,\omega_x(v)$, and there are $p!$ of them.
\end{proof}

\begin{theorem}[Leibniz rule]
  \label{thm:wedge-leibniz}
  \lean{LocallyConvex.Manifolds.altF_comp_perm, LocallyConvex.Manifolds.altF_sum_perm,
    LocallyConvex.Manifolds.alternatizeUncurryFin_eq_altF,
    LocallyConvex.Manifolds.alternatizeUncurryFin_wedge_left,
    LocallyConvex.Manifolds.alternatizeUncurryFin_wedge_right,
    LocallyConvex.Manifolds.IsLocalForm.dLoc_wedgeLoc,
    LocallyConvex.Manifolds.DifferentialForm.d_wedge_apply,
    LocallyConvex.Manifolds.DifferentialForm.d_wedge}
  \leanok
  \uses{def:wedge, thm:diff-form-d}
  For $\omega \in \Omega^p(M, E_1)$ and $\eta \in \Omega^q(M, E_2)$,
  \[
    \dd(\omega \wedge_\beta \eta) = \dd\omega \wedge_\beta \eta + (-1)^p\, \omega \wedge_\beta
    \dd\eta.
  \]
  (The degree $(p + 1) + q$ of $\dd\omega \wedge_\beta \eta$ is identified with $(p + q) + 1$.)
\end{theorem}

\begin{proof}
  \leanok
  The rule is local. By the product rule for $\beta$, the derivative of $z \mapsto
  \omega(z) \wedge_\beta \eta(z)$ at $y$ in the direction $u$ is $L(u) \wedge_\beta \eta(y) +
  \omega(y) \wedge_\beta L'(u)$, where $L(u)$, $L'(u)$ are the partial derivatives of the local
  representatives. By (3.33), it remains to prove two identities between alternating maps,
  \[
    \sum_i (-1)^i \bigl(L(v_i) \wedge B\bigr)(\dots \hat v_i \dots) = (\mathrm{Alt}L)
    \wedge B, \qquad \sum_i (-1)^i \bigl(A \wedge L'(v_i)\bigr)(\dots \hat v_i \dots) =
    (-1)^p A \wedge (\mathrm{Alt} L'),
  \]
  where $(\mathrm{Alt}L)(v_0, \dots, v_p) = \sum_i (-1)^i L(v_i)(\dots \hat v_i \dots)$. Write
  $\mathrm{alt}f(v) = \sum_\sigma \operatorname{sgn}(\sigma) f(v \circ \sigma)$ for any function $f$
  of $k$ vectors. Then $\mathrm{alt} f(v \circ \sigma) = \operatorname{sgn}(\sigma)\,
  \mathrm{alt}f(v)$, $\mathrm{alt} A = k!\, A$ for alternating $A$, and alternating a partial
  alternation $\sum_\tau \operatorname{sgn}(\tau) g(w \circ \tilde\tau)$ over a subgroup of
  order $a!$ gives $a!\, \mathrm{alt}g$. The wedge product is $(p!\,q!)^{-1}\,\mathrm{alt}$ of
  $w \mapsto \beta(f(w_1, \dots, w_p), g(w_{p+1}, \dots))$, and $\mathrm{Alt} L = (n!)^{-1}
  \mathrm{alt}\bigl(w \mapsto L(w_0)(w_1, \dots, w_n)\bigr)$, because its $i$-th term is a
  reindexing by the cycle $(0\,1 \cdots i)$, of sign $(-1)^i$. With these rules both sides of each
  identity become $(p!\,q!)^{-1}$ times the alternation of the same function, up to reindexing; in
  the second identity the reindexing is the cycle $(0\,1 \cdots p)$, which gives the sign $(-1)^p$.
\end{proof}

\section{Insertion and the Lie derivative}
\label{sec:lie-derivative}

\begin{definition}[Insertion]
  \label{def:interior}
  \lean{LocallyConvex.Manifolds.SmoothVectorField,
    LocallyConvex.Manifolds.DifferentialForm.interior,
    LocallyConvex.Manifolds.DifferentialForm.localRep_interior}
  \leanok
  \uses{def:diff-form, def:vector-field}
  For a smooth vector field $X$ and $\omega \in \Omega^{n+1}(M, F)$,
  $(i_X\omega)_x(v_1, \dots, v_n) = \omega_x(X(x), v_1, \dots, v_n)$
  \cite[Definition~3.6.15]{GN26}. In every chart, its local representative is
  $(y, v) \mapsto \omega(y)(\tilde X(y), v)$.
\end{definition}

\begin{definition}[Lie derivative]
  \label{def:lie-deriv}
  \lean{LocallyConvex.Manifolds.DifferentialForm.lieDeriv,
    LocallyConvex.Manifolds.DifferentialForm.lieDeriv_succ,
    LocallyConvex.Manifolds.DifferentialForm.lieDerivLinear}
  \leanok
  \uses{def:interior, prop:d-linear}
  We define $L_X = \dd \circ i_X + i_X \circ \dd$, with $i_X = 0$ on functions. Then $L_X$ is a
  linear map of $\Omega^n(M, F)$, and Cartan's formula \cite[Proposition~3.6.16(c)]{GN26} holds
  by definition. Theorem~\ref{thm:lie-deriv-local} shows that this agrees with the definition
  \cite[Definition~3.6.14]{GN26}.
\end{definition}

\begin{theorem}[The Lie derivative in charts]
  \label{thm:lie-deriv-local}
  \lean{LocallyConvex.Manifolds.DifferentialForm.localRep_lieDeriv,
    LocallyConvex.Manifolds.IsLocalForm.dLoc_iLoc_add,
    LocallyConvex.Manifolds.IsLocalForm.dMB_comp_vectors}
  \leanok
  \uses{def:lie-deriv, thm:diff-form-d}
  In every chart,
  \[
    (L_X\omega)(y)(v_1, \dots, v_n) = \partial_{\tilde X(y)}\omega(y)(v_1, \dots, v_n)
      + \sum_{k=1}^n \omega(y)(v_1, \dots, \dd\tilde X(y)\,v_k, \dots, v_n).
  \]
\end{theorem}

\begin{proof}
  \leanok
  Differentiating $z \mapsto \omega(z)(\tilde X(z), w)$ gives a partial derivative of $\omega$ and
  the term $\omega(y)(\dd\tilde X(y)\,u, w)$. In $\dd(i_X\omega) + i_X(\dd\omega)$ the partial
  derivatives $\partial_{v_j}\omega(y)(\tilde X(y), \dots)$ cancel, and moving
  $\dd\tilde X(y)\,v_j$ to the $j$-th place removes the signs.
\end{proof}

\begin{theorem}[Evaluation on vector fields]
  \label{thm:fields-formulas}
  \lean{LocallyConvex.Manifolds.DifferentialForm.d_apply_fields,
    LocallyConvex.Manifolds.DifferentialForm.lieDeriv_apply_fields,
    LocallyConvex.Manifolds.sum_neg_one_pow_smul_sum_update_eq}
  \leanok
  \uses{thm:lie-deriv-local, def:man-lie-bracket, thm:man-lie-bracket}
  For smooth vector fields $X_0, \dots, X_{n+1}$ and $\omega \in \Omega^{n+1}(M, F)$,
  \[
    \dd\omega(X_0, \dots, X_{n+1}) = \sum_i (-1)^i X_i.\omega(\dots, \widehat{X_i}, \dots)
      + \sum_{i<j} (-1)^{i+j} \omega([X_i, X_j], \dots, \widehat{X_i}, \dots, \widehat{X_j},
      \dots),
  \]
  \cite[Proposition~3.6.7, (3.29)]{GN26}, and $(L_Y\omega)(X_1, \dots, X_n) =
  Y.\omega(X_1, \dots, X_n) - \sum_j \omega(X_1, \dots, [Y, X_j], \dots, X_n)$
  \cite[(3.35)]{GN26}.
\end{theorem}

\begin{proof}
  \leanok
  In a chart, $X_i.\omega(\dots, \widehat{X_i}, \dots)$ is a partial derivative of $\omega$ plus
  the terms with $\dd\tilde X_k \cdot \tilde X_i$; the partial derivatives add up to
  $\dd\omega$, and the other terms regroup in pairs into the brackets. The formula for $L_Y$
  follows in the same way from Theorem~\ref{thm:lie-deriv-local}.
\end{proof}

\begin{theorem}
  \label{thm:cartan-calculus}
  \lean{LocallyConvex.Manifolds.DifferentialForm.lieDeriv_interior_sub,
    LocallyConvex.Manifolds.DifferentialForm.lieDeriv_d,
    LocallyConvex.Manifolds.DifferentialForm.lieDeriv_lieDeriv_sub,
    LocallyConvex.Manifolds.DifferentialForm.IsClosed.lieDeriv_eq,
    LocallyConvex.Manifolds.DifferentialForm.IsClosed.isExact_lieDeriv}
  \leanok
  \uses{thm:lie-deriv-local, def:man-lie-bracket, thm:diff-form-dd}
  For smooth vector fields $X$ and $Y$ \cite[Proposition~3.6.16]{GN26}:
  \begin{enumerate}
    \item[(a)] $[L_X, L_Y] = L_{[X, Y]}$;
    \item[(b)] $[L_X, i_Y] = i_{[X, Y]}$;
    \item[(d)] $L_X \circ \dd = \dd \circ L_X$;
    \item[(e)] $L_X$ maps closed forms to exact forms.
  \end{enumerate}
\end{theorem}

\begin{proof}
  \leanok
  (b) is a computation in a chart with Theorem~\ref{thm:lie-deriv-local}: the difference is
  $\omega(y)(\dd\tilde Y(y)\tilde X(y) - \dd\tilde X(y)\tilde Y(y), \dots)$, and the first
  entry is the local representative of $[X, Y]$. (d) and (e) follow from Cartan's formula and
  $\dd \circ \dd = 0$, and (a) follows algebraically from (b), (d) and Cartan's formula. Glöckner
  and Neeb deduce (a)--(d) from the corresponding identities for Lie algebra cochains.
\end{proof}

\begin{proposition}[Related vector fields]
  \label{prop:related-forms}
  \lean{LocallyConvex.Manifolds.DifferentialForm.pullback_interior,
    LocallyConvex.Manifolds.DifferentialForm.pullback_lieDeriv}
  \leanok
  \uses{def:interior, def:lie-deriv, thm:diff-form-naturality, thm:man-related}
  If $X$ and $Y$ are $\varphi$-related, then $\varphi^*(i_Y\omega) = i_X(\varphi^*\omega)$ and
  $\varphi^*(L_Y\omega) = L_X(\varphi^*\omega)$ \cite[Exercise~3.6.6]{GN26}.
\end{proposition}

\begin{proof}
  \leanok
  The first identity holds because $T\varphi$ maps $X(x)$ to $Y(\varphi(x))$; the second follows
  from it, Cartan's formula and Theorem~\ref{thm:diff-form-naturality}.
\end{proof}

\section{De Rham cohomology}

\begin{definition}[De Rham cohomology]
  \label{def:de-rham}
  \lean{LocallyConvex.Manifolds.DifferentialForm.closedForms,
    LocallyConvex.Manifolds.DifferentialForm.exactForms,
    LocallyConvex.Manifolds.DifferentialForm.exactForms_le_closedForms,
    LocallyConvex.Manifolds.DeRham}
  \leanok
  \uses{prop:d-linear, def:closed-exact}
  Let $Z^n_{dR}(M, F) = \ker \dd$ and $B^n_{dR}(M, F) = \dd(\Omega^{n-1}(M, F))$, with
  $B^0_{dR}(M, F) = 0$. Then $B^n_{dR} \subseteq Z^n_{dR}$, and the de Rham cohomology is
  $H^n_{dR}(M, F) = Z^n_{dR}(M, F) / B^n_{dR}(M, F)$ \cite[Definition~3.6.11]{GN26}.
\end{definition}

\begin{theorem}
  \label{thm:closed-locally-constant}
  \lean{LocallyConvex.Manifolds.DifferentialForm.isClosed_iff_isLocallyConstant,
    LocallyConvex.Manifolds.deRhamZeroEquiv}
  \leanok
  \uses{def:de-rham, prop:d-function}
  A function is closed if and only if it is locally constant; hence $H^0_{dR}(M, F)$ is the
  space of locally constant functions $M \to F$ \cite[Remark~3.6.12(b)]{GN26}.
\end{theorem}

\begin{proof}
  \leanok
  In a chart, a $C^1$ map with vanishing derivative is constant on convex open sets, and every
  point of a locally convex space has a convex open neighbourhood in a given open set.
\end{proof}

\begin{proposition}[Functoriality]
  \label{prop:de-rham-functorial}
  \lean{LocallyConvex.Manifolds.deRhamMap, LocallyConvex.Manifolds.deRhamMap_id,
    LocallyConvex.Manifolds.deRhamMap_comp}
  \leanok
  \uses{def:de-rham, thm:diff-form-naturality, prop:pullback-functorial}
  A smooth map $\varphi \colon N \to M$ induces linear maps $\varphi^* \colon H^n_{dR}(M, F)
  \to H^n_{dR}(N, F)$, $[\alpha] \mapsto [\varphi^*\alpha]$, with $\id^* = \id$ and
  $(\varphi \circ \psi)^* = \psi^* \circ \varphi^*$ \cite[Section~3.6]{GN26}.
\end{proposition}

\begin{proof}
  \leanok
  By Theorem~\ref{thm:diff-form-naturality}, pullbacks of closed (exact) forms are closed
  (exact); functoriality comes from Proposition~\ref{prop:pullback-functorial}.
\end{proof}

\section{The Poincaré Lemma}

Homotopies are smooth maps $H \colon \R \times M \to N$, with $H_t = H(t, \cdot)$; on $\R \times M$
we use the product charts $\id_{\R} \times \varphi$. Glöckner and Neeb use $[0, 1] \times M$, a
manifold with boundary; a homotopy on $[0, 1] \times M$ can be precomposed with a smooth map
$\R \to [0, 1]$ which is $0$ near $(-\infty, 0]$ and $1$ near $[1, \infty)$.

\begin{definition}[Fiber integral]
  \label{def:fiber-integral}
  \lean{LocallyConvex.Manifolds.fibLoc, LocallyConvex.Manifolds.IsLocalForm.isLocalForm_fibLoc,
    LocallyConvex.Manifolds.DifferentialForm.fibInt,
    LocallyConvex.Manifolds.DifferentialForm.localRep_fibInt}
  \leanok
  \uses{def:diff-form}
  Let $F$ be Mackey complete. For $\omega \in \Omega^{k+1}(\R \times M, F)$, the fiber integral
  $I(\omega) \in \Omega^k(M, F)$ is
  \[
    I(\omega)_x(v_1, \dots, v_k) = \int_0^1 \omega_{(t, x)}\bigl((1, 0), (0, v_1), \dots,
    (0, v_k)\bigr)\, dt
  \]
  (weak integral; in a chart, the integrand is a smooth curve, whose weak integral exists because
  $F$ is Mackey complete). In the chart $\varphi$, the local representative of $I(\omega)$ is the
  same integral of the local representative of $\omega$ in the chart $\id \times \varphi$, so
  $I(\omega)$ is smooth by the smoothness of parameter-dependent integrals.
\end{definition}

\begin{theorem}[Homotopy formula {\cite[(3.40)]{GN26}}]
  \label{thm:homotopy-formula}
  \lean{LocallyConvex.Manifolds.IsLocalForm.dMB_fibLoc,
    LocallyConvex.Manifolds.IsLocalForm.dLoc_fibLoc_add,
    LocallyConvex.Manifolds.IsLocalForm.fibLoc_dLoc_zero,
    LocallyConvex.Manifolds.DifferentialForm.pullback_sliceIncl_apply,
    LocallyConvex.Manifolds.DifferentialForm.d_fibInt_add,
    LocallyConvex.Manifolds.DifferentialForm.fibInt_d_zero}
  \leanok
  \uses{def:fiber-integral, thm:diff-form-d, def:diff-form-pullback}
  With $i_t(x) = (t, x)$, $\dd I(\omega) + I(\dd\omega) = i_1^*\omega - i_0^*\omega$ for
  $\omega \in \Omega^{k+1}(\R \times M, F)$, and $I(\dd f) = i_1^* f - i_0^* f$ for functions $f$.
\end{theorem}

\begin{proof}
  \leanok
  Glöckner and Neeb use Cartan's formula. We compute in a chart. Differentiating under the
  integral sign (which requires $\Ck{2}$ integrands in Mackey complete spaces),
  $\dd I(\omega)(v_0, \dots, v_k) = \sum_j (-1)^j \int_0^1 \partial_{(0, v_j)}
  \omega_{(t,x)}((1,0), (0, v_0), \dots, \widehat{(0, v_j)}, \dots)\, dt$. In the formula (3.33)
  for $(\dd\omega)_{(t,x)}((1, 0), (0, v_0), \dots)$, the terms with $j \geq 1$ are the same with
  the opposite sign, and the remaining term is $\partial_t \omega_{(t,x)}((0, v_0), \dots)$. By the
  fundamental theorem of calculus, its integral is $\omega_{(1, x)}(\dots) - \omega_{(0,
  x)}(\dots) = (i_1^*\omega - i_0^*\omega)_x(v)$.
\end{proof}

\begin{theorem}[Homotopy Lemma {\cite[Lemma~3.6.25]{GN26}}]
  \label{thm:homotopy-lemma}
  \lean{LocallyConvex.Manifolds.DifferentialForm.pullback_sub_pullback_eq_d_fibInt,
    LocallyConvex.Manifolds.DifferentialForm.pullback_eq_pullback_of_isClosed_zero,
    LocallyConvex.Manifolds.deRhamMap_homotopy}
  \leanok
  \uses{thm:homotopy-formula, prop:de-rham-functorial, prop:pullback-functorial}
  Let $H \colon \R \times M \to N$ be smooth and $F$ Mackey complete. For a closed $\alpha \in
  \Omega^{k+1}(N, F)$, $H_1^*\alpha - H_0^*\alpha = \dd I(H^*\alpha)$, and $H_1^* f = H_0^* f$ for
  closed functions $f$. Hence $H_0^* = H_1^* \colon H^n_{dR}(N, F) \to H^n_{dR}(M, F)$ for all $n$.
\end{theorem}

\begin{proof}
  \leanok
  Apply the homotopy formula to $\omega = H^*\alpha$, for which $\dd\omega = H^*\dd\alpha = 0$, and
  use $i_t^* H^* = (H \circ i_t)^* = H_t^*$.
\end{proof}

\begin{theorem}[Poincaré Lemma {\cite[Definition~3.6.27, Theorem~3.6.28]{GN26}}]
  \label{thm:poincare}
  \lean{LocallyConvex.Manifolds.SmoothlyContractible,
    LocallyConvex.Manifolds.DifferentialForm.pullback_const_eq_zero,
    LocallyConvex.Manifolds.DifferentialForm.IsClosed.isExact_of_smoothlyContractible,
    LocallyConvex.Manifolds.deRham_subsingleton_of_smoothlyContractible}
  \leanok
  \uses{thm:homotopy-lemma}
  Let $M$ be smoothly contractible: there are $p \in M$ and a smooth $H \colon \R \times M \to M$
  with $H(0, x) = x$ and $H(1, x) = p$. Then every closed form of positive degree on $M$ with
  values in a Mackey complete space is exact, i.e.\ $H^n_{dR}(M, F) = 0$ for $n \geq 1$.
\end{theorem}

\begin{proof}
  \leanok
  $H_0 = \id_M$ and $H_1$ is constant, so $H_1^*\alpha = 0$ for $\deg \alpha \geq 1$. By the
  Homotopy Lemma, $-\alpha = \dd I(H^*\alpha)$.
\end{proof}

\begin{corollary}[{\cite[Corollary~3.6.29]{GN26}}]
  \label{cor:poincare-star}
  \lean{LocallyConvex.Manifolds.mbGroupoid, LocallyConvex.Manifolds.smoothlyContractible_of_starConvex,
    LocallyConvex.Manifolds.DifferentialForm.IsClosed.isExact_of_starConvex}
  \leanok
  \uses{thm:poincare}
  Let $U \subseteq E$ be open and star-shaped with respect to $p \in U$. Every closed form of
  positive degree on $U$ with values in a Mackey complete space is exact.
\end{corollary}

\begin{proof}
  \leanok
  Open subsets of manifolds are manifolds, because restrictions of smooth diffeomorphisms to open
  sets are smooth diffeomorphisms. With $c(t) = (1 - t)^2/((1 - t)^2 + t^2)$, which is smooth with
  values in $[0, 1]$, $c(0) = 1$ and $c(1) = 0$, the map $H(t, x) = p + c(t)(x - p)$ contracts $U$
  to $p$.
\end{proof}

\begin{corollary}[{\cite[Corollary~3.6.26]{GN26}}]
  \label{cor:de-rham-product}
  \lean{LocallyConvex.Manifolds.smoothMBM_scaleFiber,
    LocallyConvex.Manifolds.smoothMBM_zeroSection,
    LocallyConvex.Manifolds.deRhamProdEquiv, LocallyConvex.Manifolds.coe_deRhamProdEquiv}
  \leanok
  \uses{thm:homotopy-lemma, prop:de-rham-functorial}
  Let $G$ be a locally convex space and $F$ a Mackey complete one. The projection
  $p \colon M \times G \to M$ induces isomorphisms
  $p^* \colon H^n_{dR}(M, F) \to H^n_{dR}(M \times G, F)$, with inverse $i_0^*$, where
  $i_0(m) = (m, 0)$.
\end{corollary}

\begin{proof}
  \leanok
  $p \circ i_0 = \id_M$, and $H(t, (m, x)) = (m, t x)$ is a smooth homotopy from $i_0 \circ p$
  to $\id_{M \times G}$, so $p^* i_0^* = (i_0 \circ p)^* = \id$ by the Homotopy Lemma.
\end{proof}

\begin{theorem}[The Lie derivative along a local flow {\cite[Remark~3.6.17, (3.36)]{GN26}}]
  \label{thm:lie-flow}
  \lean{LocallyConvex.Manifolds.IsLocalForm.hasDirDerivAt_flow,
    LocallyConvex.Manifolds.IsLocalFlow,
    LocallyConvex.Manifolds.IsLocalFlow.hasDirDerivAt_pullback}
  \leanok
  \uses{thm:lie-deriv-local}
  Let $X$ be a smooth vector field with a local $C^2$-flow $\mathrm{Fl}$. Then
  $\frac{d}{dt}\big|_{t=0} ((\mathrm{Fl}_t)^*\omega)_x(v_1, \dots, v_n) = (L_X\omega)_x(v_1, \dots,
  v_n)$.
\end{theorem}

\begin{proof}
  \leanok
  In a chart, differentiate $t \mapsto \omega(\varphi_t(y))(\dd\varphi_t(y)v_1, \dots)$, where
  $\varphi$ is the local flow: the base point moves with velocity $X(y)$, and
  $\partial_t \dd\varphi_t(y) v = \dd X(y) v$ at $t = 0$ by exchanging the order of
  differentiation in $\partial_t \varphi = X \circ \varphi$. This is the local formula for $L_X$.
  Glöckner and Neeb assume $C^1$-flows; we do not claim this argument is new.
\end{proof}
